\section{Constructing the original patched module and all support completions}
\label{sec:patched-support-v040}

\status{PMS constructs the patched module from the original finite-level
modules, proves the complete original S5 support theorem, and supplies its
finite-level arithmetic hypotheses on the same Lim, TWAut and universal
ring maps.  The ordinary automorphic producer remains an input.  All
constructions retain the original ideal, coefficient action and
ultrafilter.}

The source pin is
\par\smallskip\noindent
\texttt{7cbf7a0c32fca1badf05931879a37c3f8a2b03b4}.
\par\smallskip\noindent
Within \texttt{AutomorphicVarieties}, write $\mathcal P$ for
\texttt{PatchingModules/}.
\par\noindent
Write $\mathcal H$ for its \texttt{AnnihilatorBaseChange/}
subdirectory.  Formal declarations below
identify exact contracts; their invocation is not the proof.  Standard
Noetherian completion facts are the checked applications recorded with
this chapter.  The actual inverse limit, module, fiber and dimension
arguments are given here.

\subsection{One ultrafilter and an actual finite patched module}

% BEGIN PMS1
\paragraph{PMS1. The complete input contract and its quantifier order.}
\label{claim:PMS1}
Let $L_0$ be a complete Noetherian local ring with finite residue field,
$R$ a Noetherian local ring, $\sigma$ a finite set, and
$s:L_0[[Y_a:a\in\sigma]]\to R$ a ring map.  Put
\[
 c=s|_{L_0},\quad y_a=s(Y_a),\quad I=(y_a),\quad
 J_j=\mathfrak m_{L_0}^{,j}R,\quad K_{k,j}=I^k+J_j.
                                                        \tag{PMS.1}
\]
Here $J_j$ is formed using $c$.  It is not $\mathfrak m_R^j$.
Fix one ultrafilter $\mathscr U$ on $\mathbb N$ containing every tail,
one $f\in R$, an ideal $\mathfrak a\subset R$, and actual $R$-modules
$M_N$.  Let $X$ be a finite $L_0$-module with a commuting $R$-action
satisfying $a x=c(a)x$ and $y_aX=0$.  Assume:
\begin{enumerate}
\item For every $k,j$, on a $\mathscr U$-large set of levels there are
an integer $r_N$ and an additive isomorphism
\[
 M_N/K_{k,j}M_N\simeq
 \bigl(L_0[[Y]]/((Y)^k+\mathfrak m_{L_0}^{,j})\bigr)^{r_N}
                                                        \tag{PMS.2}
\]
intertwining multiplication by every power series through $s$.
\item For every $j$, on a $\mathscr U$-large set there are $R$-linear
maps in both directions between $X/K_{1,j}X$ and $M_N/K_{1,j}M_N$;
both composites are multiplication by this same $f$.
\item There are one integer $e\geq0$ and one set $G_s\subset R$
with $(G_s)=\mathfrak a$ such that, for every $r\in G_s$ and $j$,
on a $\mathscr U$-large set there is $b\in R$ with
\[
 (fr)^e-b\in\mathfrak m_R^j,\qquad bM_N=0.        \tag{PMS.3}
\]
\end{enumerate}
The events and ranks may depend on the displayed indices.  The exponent
$e$ is chosen before $r,j,N$.  No bounded rank, common basis at all
levels, compatible choice of the finite-level maps, or interchange of
these quantifiers is assumed.  These are the literal three antecedents
of \texttt{S5}, except that freeness of $X$ is needed only from PMS14
onward.  The generic module construction therefore also covers the
finite-$X$ scope of \texttt{pmodEx}.

All rings used below remain Noetherian: localization preserves that
property; completion at a finitely generated ideal preserves it.
Completion is exact on finite modules and equals tensoring with the
completed ring.  Its quotients by powers of the completion ideal are
the original quotients.  These facts apply to $R_f$, its ideal $IR_f$,
and to the actual localizations at specified primes.  The precise
primary statements are Stacks tags
\href{https://stacks.math.columbia.edu/tag/00MA}{00MA},
\href{https://stacks.math.columbia.edu/tag/00MB}{00MB},
\href{https://stacks.math.columbia.edu/tag/031C}{031C} and
\href{https://stacks.math.columbia.edu/tag/05GH}{05GH}; the needed
Noetherian and finite-generation hypotheses have just been specified.
\hfill\textit{Source: \texttt{PatchedSupportCompletions.S5}.}
% END PMS1

% BEGIN PMS2
\paragraph{PMS2. A diagonal ideal criterion for the actual ultralimit.}
\label{claim:PMS2}
For an antitone sequence of ideals $H_j\subset R$ define
\[
 U(H,M)=\prod_N M_N\Big/
 \{p:\ \forall j,\ p_N\in H_jM_N\text{ for }\mathscr U
       \text{-almost every }N\}.                  \tag{PMS.4}
\]
This is the source $\mathrm{ULim}$, not a limit of selected points.
For every finitely generated ideal $A\subset R$,
\[
 [p]\in A U(H,M)
 \quad\Longleftrightarrow\quad
 \forall j,\ p_N\in(A+H_j)M_N
                    \text{ for }\mathscr U\text{-almost every }N.
                                                        \tag{PMS.5}
\]
Choose generators $a_1,\ldots,a_v$ of $A$.  The forward implication
follows by lifting the finitely many coefficient sequences in
$[p]=\sum a_i[q_i]$.  For the reverse implication let $E_j$ be the
large event at precision $j$, replacing it by
$\{N\geq j\}\cap\bigcap_{i\leq j}E_i$.  These events are decreasing.
For each $N$ let $\nu(N)$ be the largest $j\leq N$ for which
$N\in E_j$, and set it to zero when there is none.  Then
$\nu\to\infty$ along $\mathscr U$, and on a large set
$p_N=\sum_i a_iq_{i,N}+z_N$ with $z_N\in H_{\nu(N)}M_N$.
Set unused coefficients to zero off that event.  For each fixed $j$,
$z_N\in H_jM_N$ eventually along $\mathscr U$; hence $[z]=0$.
This proves (PMS.5), without countable completeness of the ultrafilter.

In particular, if $B\subset A$ are finitely generated ideals, the
identity on products induces a surjection
$U(B+H_j,M)\to U(A+H_j,M)$ with kernel $A$ times the source.
The target is killed by $A$: this is immediate on representatives.
Surjectivity uses the same representatives, while the kernel assertion
is exactly (PMS.5).  The diagonal construction will also be used to
choose finite-level maps at a precision tending to infinity.
\hfill\textit{Sources: \texttt{UltraDefs}, \texttt{ULimKer},
\texttt{ULimTrans}.}
% END PMS2

% BEGIN PMS3
\paragraph{PMS3. The finite special fiber determines the first ultralimit.}
\label{claim:PMS3}
For the constant system $X$ the diagonal map is an isomorphism
\[
 X\ \simeq\ U(K_{1,j},X).                         \tag{PMS.6}
\]
Indeed $K_{1,j}X=\mathfrak m_{L_0}^{,j}X$.  Each quotient is a
finite set: the graded pieces of $L_0/\mathfrak m_{L_0}^j$ are finite
dimensional over its finite residue field, and $X$ is finite over
$L_0$.  A sequence in a finite set has one unique value on a
$\mathscr U$-large set.  These values for the quotients are compatible;
completeness of the finite $L_0$-module $X$ supplies their unique lift.
Separation proves injectivity.  The isomorphism is $R$-linear because
each $r\in R$ acts $L_0$-linearly and preserves every
$\mathfrak m_{L_0}^jX$.

Set $U_1=U(K_{1,j},M)$.  Apply the diagonal choice of PMS2 to the
existence events for the paired maps in PMS1.  At level $N$ choose a
pair at precision $\nu(N)$ and arbitrary representatives of its
values.  These give maps on the product quotients: additivity,
$R$-linearity, independence of representatives and preservation of
the null sequences all hold modulo $K_{1,\nu(N)}$, and therefore
hold in the ultralimit.  For preservation at a fixed precision $j$,
intersect with $\nu(N)\geq j$ and use the linearity of the chosen map
modulo $K_{1,\nu(N)}$.  Its image of $K_{1,j}$ still lies in
$K_{1,j}$ there.  Thus (PMS.6) gives actual $R$-linear maps
\[
 \Phi:X\longrightarrow U_1,\qquad
 \Psi:U_1\longrightarrow X,
 \quad \Phi\Psi=f,\quad\Psi\Phi=f.               \tag{PMS.7}
\]
The composite identities follow at each chosen precision and then at
every fixed precision.  No compatibility between different chosen
levels was used.  On $S=R_f$, the map $\Psi$ becomes an isomorphism
$S\otimes_R U_1\simeq S\otimes_R X$, with inverse $f^{-1}\Phi$.
\hfill\textit{Source: \texttt{ExistsUltraLimitPairOfEventualQuotientMaps}.}
% END PMS3

% BEGIN PMS4
\paragraph{PMS4. Truncated freeness gives homogeneous independence.}
\label{claim:PMS4}
For a module $V$ and an auxiliary ideal $J$, write
$\mathrm{QR}_a(V;y,J)$ for the assertion
\[
 \sum_{|\alpha|=a}y^\alpha v_\alpha
       \in(I^{a+1}+J)V
 \quad\Longrightarrow\quad
 v_\alpha\in(I+J)V\quad\hbox{for every }\alpha.  \tag{PMS.8}
\]
Here the sum uses all multiindices of degree $a$ on the finite set
$\sigma$.  In a formal power-series ring modulo
$(Y)^{a+1}+\mathfrak m_{L_0}^j$, the degree-$a$ coefficient of
$\sum Y^\alpha F_\alpha$ is the constant coefficient of the
corresponding $F_\alpha$, reduced modulo $\mathfrak m_{L_0}^j$.
Consequently (PMS.2), with $k=a+1$, implies
$\mathrm{QR}_a(M_N;y,J_j)$ on a large set.  To spell out the return
from coordinates, lift a basis of the finite free quotient, write
the coordinate vector of $v_\alpha$, and subtract its expression
in these basis lifts.  Its coordinate constants vanish modulo
$\mathfrak m_{L_0}^j$, so the expression lies in $(I+J_j)M_N$;
the error lies in $K_{a+1,j}M_N\subset K_{1,j}M_N$.
This is valid also for rank zero and for $a=0$.

Put $U_k=U(K_{k,j},M)$.  For $a<k$, the same property gives
$\mathrm{QR}_a(U_k;y,0)$.  In fact, lift the finitely many
$v_\alpha$ to sequences.  By PMS2, a relation modulo $I^{a+1}U_k$
holds at every precision modulo $I^{a+1}+J_j$, because
$I^k\subset I^{a+1}$.  Apply the level property and then PMS2
again with $A=I$ to each coefficient.  Thus the full homogeneous
independence needed at every finite degree survives the one
ultrafilter quotient.
\hfill\textit{Sources: \texttt{QuotientRegularAtOfQuotientFreeOverPowerSeries},
\texttt{QRegLim}.}
% END PMS4

% BEGIN PMS5
\paragraph{PMS5. Localize the complete quotient tower.}
\label{claim:PMS5}
Let $S=R_f$, $I_S=IS$, and $P_k=S\otimes_R U_k$.
Localization is exact, so PMS2 gives surjective transition maps
\[
 P_{k+1}\longrightarrow P_k,
 \quad\ker=I_S^kP_{k+1},\quad I_S^kP_k=0.         \tag{PMS.9}
\]
In particular $P_0=0$.  PMS3 identifies $P_1$ with $S\otimes_R X$,
which is finite over $S$: any finite $L_0$-generating list of $X$
also generates it over $R$.
Let
\[
 B_B=\widehat{S}_{I_S}=\varprojlim_k S/I_S^k.      \tag{PMS.10}
\]
Every $P_k$ becomes a $B_B$-module through
$B_B/I_S^kB_B=S/I_S^k$, and the transitions are $B_B$-linear.
All modules and actions are those just constructed.  Quasiregularity
in degrees $a<k$ persists under localization: clearing the finitely
many denominators in a relation reduces it to PMS4, and clearing the
denominators in the conclusions gives (PMS.8) over $S$.

The ring (PMS.10) has every field of the source $\mathrm{IsYCpl}(y,f)$.
It is Noetherian, and $f$ is a unit.  For $z\in B_B$ and $k$, choose
$r/f^v\in S$ with the same image modulo $I_S^k$.  Then
$f^vz-r\in I^kB_B$, proving the required density with denominators.
Conversely, if $r\in R$ maps into $I^kB_B$, the unchanged quotient
$S/I_S^k$ gives $r/1\in I_S^k$; clearing a denominator yields
$f^vr\in I^k$.  This is the exact strictness field, including $k=0$.
\hfill\textit{Sources: \texttt{Patched}, \texttt{TowerCompl},
\texttt{YCplMemOf.yCpl}.}
% END PMS5

% BEGIN PMS6
\paragraph{PMS6. Construct a finite module with every prescribed quotient.}
\label{claim:PMS6}
We prove the required tower statement, rather than assuming that its
inverse limit is finite.  Let $B$ be complete for a finitely generated
ideal $A$, and suppose a tower $P_k$ satisfies (PMS.9), with $I_S$
replaced by $A$, and $P_1$ finite.  Choose generators of $P_1$ and
lift them successively along the surjections.  They give compatible
maps $\psi_k:B^v\to P_k$.  Each map is onto: its image surjects
modulo $A$, and repeated substitution in
$P_k=\operatorname{im}\psi_k+AP_k$, using $A^kP_k=0$, removes
the final error.  This is nilpotent Nakayama and does not presume
that $P_k$ was already finite.

Set $K_k=\ker\psi_k$.  Then
\[
 A^kB^v\subset K_k,\quad K_{k+1}\subset K_k,
 \quad K_k\subset A^kB^v+K_{k+1}.                \tag{PMS.11}
\]
For the last inclusion, the image in $P_{k+1}$ of an element of
$K_k$ lies in $A^kP_{k+1}$; use the surjectivity of $\psi_{k+1}$
to lift its finitely many coefficients.

Let $M=\varprojlim P_k$ as a submodule of the product, with its
literal compatible-coordinate action.  For $z\in M$, choose a
coefficient vector correct in $P_1$.  If it is correct in $P_k$,
(PMS.11) permits a correction in $A^kB^v$ that makes it correct
in $P_{k+1}$.  The resulting coefficient vectors converge in
$B^v$.  Their limit maps to $z$ in every $P_k$, since
$A^kB^v\subset K_k$.  Thus $B^v\to M$ is onto, and $M$ is finite.
Every coordinate map $q_k:M\to P_k$ is onto as well.

Its kernel is exactly $A^kM$.  One inclusion is immediate.  For
the other, lift $z\in\ker q_k$ to $b\in K_k$.  Repeatedly use
(PMS.11) to write
$b=d_k+d_{k+1}+\cdots+d_{r-1}+b_r$, where
$d_j\in A^jB^v$ and $b_r\in K_r$.
The convergent sum $d=\sum_{j\geq k}d_j$ lies in $A^kB^v$.
For every $r$, $b-d\in K_r$: the finite error is in $K_r$ and
the tail is in $A^rB^v\subset K_r$.
Hence $z$ is the image of $d$ and belongs to $A^kM$.
Finally $\bigcap_k\ker q_k=0$ because $M$ is a submodule of the
product.  Applying this proof to $B=B_B$, $A=IB_B$ defines the
actual patched module $M_M$ and proves
\[
 M_M\text{ finite over }B_B,\qquad
 M_M/I^kM_M\simeq P_k\quad(k\geq0),\qquad
 \bigcap_k I^kM_M=0.                             \tag{PMS.12}
\]
\hfill\textit{Sources: \texttt{TowerFin.exists\_tower\_limit},
\texttt{TowerCompl}.}
% END PMS6

% BEGIN PMS7
\paragraph{PMS7. Quasiregularity and the exact scalar-compatible fiber.}
\label{claim:PMS7}
For every $a$, apply $q_{a+1}$ to a degree-$a$ relation in $M_M$.
By PMS4--PMS5 its coefficients map into $IP_{a+1}$.  Lift these
expressions along $q_{a+1}$; their errors belong to
$I^{a+1}M_M\subset IM_M$.  Thus each coefficient is in $IM_M$.
This proves the literal $\mathrm{IsQuasiRegularSeq}(M_M,y)$.

The isomorphism $P_1\simeq S\otimes_R X$ of PMS3 is $S$-linear.
Since $IX=0$ and $B_B/IB_B=S/I_S$, base change gives a natural
isomorphism
$S\otimes_R X\simeq B_B\otimes_R X$.
It follows from (PMS.12) that
\[
 M_M/IM_M\simeq_{B_B} B_B\otimes_R X.            \tag{PMS.13}
\]
The subscript is justified explicitly: a scalar $b\in B_B$ can
be represented modulo $IB_B$ by a scalar in $S$; both modules in
(PMS.13) are killed by $I$, so the $S$-linear isomorphism respects
that scalar, independently of the representative.  Thus this is
the source's entire $\mathrm{fib}$ field, not merely an additive
or coefficient-linear identification.

We shall use one annihilator consequence of the same construction.
If $r\in R$ satisfies
\[
 \forall k\ \exists v\ \forall j\ 
 f^v rM_N\subset K_{k,j}M_N
       \text{ for }\mathscr U\text{-almost every }N,             \tag{PMS.14}
\]
then $rM_M=0$.  Indeed $f^vr$ kills $U_k$, hence $r$ kills $P_k$
after localization.  The coordinate maps and their common zero
kernel in (PMS.12) give the conclusion.  The exponent $v$ may
depend on $k$, but is before $j$; that order will be checked below.
\hfill\textit{Sources: \texttt{QRegTower}, \texttt{FibEquiv},
\texttt{Patched}.}
% END PMS7

% BEGIN PMS8
\paragraph{PMS8. Convert approximation in the maximal ideal into annihilation.}
\label{claim:PMS8}
For fixed $j$, the finite set $X/\mathfrak m_{L_0}^jX$ is an
$R$-module.  Its sequence of submodules
$\mathfrak m_R^c(X/\mathfrak m_{L_0}^jX)$ stabilizes.  The stable
submodule is finitely generated, because it is a finite set;
Nakayama makes it zero.  Choose $c=c(j)$ killing this quotient.
On the large event of paired maps at precision $j$, their
linearity and composite identity show
\[
 f\mathfrak m_R^c M_N\subset (I+J_j)M_N.
\]
Modulo $J_jM_N$, iterate this inclusion $k$ times to obtain
\[
 f^k\mathfrak m_R^{ck}M_N\subset(I^k+J_j)M_N.    \tag{PMS.15}
\]
The case $k=0$ is immediate.  This proves the source's full
finite-quotient estimate:
for each $k,j$ there is an exponent before the large event and
before all elements of that power of $\mathfrak m_R$.

Fix $r\in G_s$.  Given $k,j$, use (PMS.3) at precision $ck$ and
intersect with the event in (PMS.15).  The annihilating $b$ there
gives
$f^k(fr)^eM_N\subset K_{k,j}M_N$.
Thus (PMS.14) holds for $(fr)^e$, with exponent $v=k$ independent
of $j$.  We have $(fr)^eM_M=0$; $f$ is a unit in $B_B$, so
$r^eM_M=0$.  For $e=0$ this says $M_M=0$, which is also allowed.
It follows for every generator, and then for the ideal they
generate, that
\[
 \mathfrak a\subset
 R\cap\sqrt{\operatorname{Ann}_{B_B}M_M}.        \tag{PMS.16}
\]
Together PMS5--PMS8 construct every field of the original
$\mathrm{PMod}(y,f,\mathfrak a,X)$: the actual ring, its strict
completion description, a finite module, quasiregularity, the
$B_B$-linear fiber, and the radical annihilator inclusion.
\hfill\textit{Sources: \texttt{ExistsPowMaximalIdealSmulFiniteEqZero},
\texttt{PmodEx.pmodEx}.}
% END PMS8

\subsection{The full support theorem, including movement between completions}

% BEGIN PMS9
\paragraph{PMS9. Compare primes and completions before and after patching.}
\label{claim:PMS9}
Let $q\subset R$ be prime, $I\subset q$, and $f\notin q$.
There is a natural identification
\[
 B_B/qB_B\simeq S/qS\simeq (R/q)_f.             \tag{PMS.17}
\]
Indeed $IB_B\subset qB_B$ and $B_B/IB_B=S/I_S$;
quotienting by the image of $q$ proves the first equality.
Consequently $qB_B$ is prime and contracts to $q$.
More generally primes of $B_B$ containing $I$ correspond to
primes of $R$ containing $I$ and avoiding $f$.  Hence, if $Q$
contains $I$ and $Q\cap R\subset q$, then $Q\subset qB_B$.

For every $v\geq1$,
$B_B/q^vB_B=S/q^vS$ as well.  Completion is exact on the finite
$S$-module $S/q^vS$, and this module is already $I_S$-adically
complete since $I_S^v$ kills it.  Localizing at the matched
primes and taking the inverse limit of these quotients gives
\[
 C_q:=\widehat{(B_B)_{qB_B}}
       \simeq\widehat{R_q}.                      \tag{PMS.18}
\]
All maps from $R$ agree: $f$ was already a unit in $R_q$.
The map $B_B\to C_q$ is flat; its localization at $qB_B$ is
faithfully flat.  Here the completion ideal is the maximal ideal
of the Noetherian local ring $(B_B)_{qB_B}$, so Stacks
\href{https://stacks.math.columbia.edu/tag/00MC}{00MC} applies.
No assertion that $B_B$ itself is local is needed.
\hfill\textit{Source: \texttt{YPrime.yPrime}.}
% END PMS9

% BEGIN PMS10
\paragraph{PMS10. Support is detected by the exact special fiber.}
\label{claim:PMS10}
For a finite module $V$ over a ring $A$ and a prime $p$,
\[
 \operatorname{Ann}_A V\subset p
 \ \Longleftrightarrow\ V_p\ne0
 \ \Longleftrightarrow\ V\otimes_A\kappa(p)\ne0.                \tag{PMS.19}
\]
For the first equivalence, if $V_p=0$, a product of denominators
killing a finite generating list kills $V$ and is outside $p$.
The second is local Nakayama.  Thus for any $A$-algebra $C$ and
prime $Q\subset C$, support of $C\otimes_A V$ at $Q$ is exactly
support of $V$ at $Q\cap A$: extension of a vector space along
the field inclusion $\kappa(Q\cap A)\to\kappa(Q)$ preserves
nonzeroness.  No flatness is required for this assertion.

Let $C$ be any $B_B$-algebra and set $V=C\otimes_{B_B}M_M$.
For a prime $Q\subset C$ containing all $y_a$, quotienting $V$
by $IV$ does not alter its residue fiber at $Q$.  By (PMS.13),
$V/IV=C\otimes_R X$.  Apply (PMS.19), first to $M_M$ or $V$
and then to the finite $R$-module $X$, to conclude
\[
 \operatorname{Ann}_{C}V\subset Q
 \quad\Longleftrightarrow\quad
 \operatorname{Ann}_{R}X\subset Q\cap R.          \tag{PMS.20}
\]
This is the exact \texttt{pmFib} equivalence.  It concerns the
annihilator of the original fiber $X$, with its original action.
\hfill\textit{Source: \texttt{PmFib.pmFib}.}
% END PMS10

% BEGIN PMS11
\paragraph{PMS11. Move support with exactly the same contracted prime.}
\label{claim:PMS11}
Let $q_i,q_{i'}$ be primes as in PMS9 and put
$C_i=C_{q_i}$, $V_i=C_i\otimes_{B_B}M_M$.
Suppose $Z\subset Q$ are primes of $C_i$,
$\operatorname{Ann}_{C_i}V_i\subset Z$, all $y_a\in Q$, and
$Q\cap R\subset q_{i'}$.
Set $p=Z\cap B_B$.  Formula (PMS.19) gives
$\operatorname{Ann}_{B_B}M_M\subset p$.
The prime $Q\cap B_B$ contains $I$, so PMS9 gives
$Q\cap B_B\subset q_{i'}B_B$; hence $p\subset q_{i'}B_B$.
Localize $p$ in $(B_B)_{q_{i'}B_B}$ and use faithful flatness of
its completion to obtain a prime $Z'$ of $C_{i'}$ contracting to
that localized prime.  To justify this last use directly, a
faithfully flat map has nonzero base change to the residue field
at any prime; a maximal ideal of that nonzero algebra provides
a prime above it.  Formula (PMS.19) then puts $Z'$ in the support
of $V_{i'}$.

The contractions of $Z'$ and $Z$ to $B_B$, and therefore to $R$,
are equal to the same $p$ and its contraction.  This proves the
full movement field, for arbitrary $i,i'$, not only movement of
minimal primes or of a chosen point.  The auxiliary prime $Q$
is used only to put $p$ inside the new completion center.
\hfill\textit{Source: \texttt{PmMove.pmMove}.}
% END PMS11

% BEGIN PMS12
\paragraph{PMS12. The completed original fiber is finite free over the coefficient completion.}
\label{claim:PMS12}
Now suppose $X$ is finite free over $L_0$.  Fix $q$ as in PMS9,
$p=q\cap L_0$, $L_a=\widehat{(L_0)_p}$ and $C=C_q$.
The local map $(L_0)_p\to R_q$ extends uniquely to a local map
$L_a\to C$ by completion.  We claim
\[
 C\otimes_R X\text{ is finite free over this }L_a.             \tag{PMS.21}
\]
The assertion includes the zero module.
If $\operatorname{Ann}_R X\not\subset q$, then $X_q=0$ by
(PMS.19), and (PMS.21) is clear.  Otherwise put
$E=R/\operatorname{Ann}_R X$.
Its faithful action embeds $E$ as an $L_0$-submodule of
$\operatorname{End}_{L_0}X$, a finite free module.  Since $L_0$
is Noetherian, $E$ is finite over $L_0$.

Localize at $p$ and complete the coefficient ring.  The finite
$L_a$-algebra $E_p\otimes_{(L_0)_p}L_a$ is a product of complete
local rings, one for each maximal ideal of $E_p$ over $p$.
This is exactly Stacks
\href{https://stacks.math.columbia.edu/tag/0325}{0325}, applied
to the Noetherian complete local ring $L_a$ and a finite algebra.
The factors are indexed as asserted because reduction modulo
$pL_a$ is the original finite-dimensional algebra $E_p/pE_p$;
its distinct maximal ideals determine its Artinian local factors.
The maximal ideal induced by $q$ selects one idempotent $e_q$.

There is a natural identification
\[
 C\otimes_R X\simeq
 e_q\bigl(L_a\otimes_{L_0}X\bigr).               \tag{PMS.22}
\]
Here are the completion details.  On $E_q$, the radical of
$pE_q$ is its maximal ideal, because $E_p/pE_p$ is Artinian and
localization selects its $q$-factor.  A power of that maximal
ideal lies in $pE_q$, so the two adic topologies are cofinal.
For each finite quotient the idempotent decomposition of the
finite algebra selects exactly the localization at $q$.
Taking inverse limits therefore identifies the right side of
(PMS.22) with the completion of $X_q$.
The latter is $\widehat{R_q}\otimes_R X$ by exactness and the
finite-module completion formula, and (PMS.18) identifies it with
the left side.  These maps agree on both the original $R$-action
and the coefficient action; no newly chosen coefficient embedding
is involved.

The right side is a direct summand of the finite free $L_a$-module
$L_a\otimes_{L_0}X$.  It is finite projective, hence free because
$L_a$ is local.  One can check the last assertion by lifting a
basis modulo the maximal ideal, applying Nakayama for generation,
and splitting the resulting surjection; its kernel has zero
residue quotient and is therefore zero.  This proves (PMS.21).
\hfill\textit{Sources: \texttt{CaCplModfree},
\texttt{FlatnessPrecompleteProducts}; used in \texttt{PmDimQuot}.}
% END PMS12

% BEGIN PMS13
\paragraph{PMS13. Quasiregularity gives a power-series basis on the actual completion.}
\label{claim:PMS13}
Let $N=C\otimes_{B_B}M_M$, with $C,L_a$ as in PMS12.
It is finite over the complete Noetherian local ring $C$.
The images of $y_a$ lie in its maximal ideal.
The sequence $y$ is quasiregular on $N$: equivalently the
canonical homogeneous maps
\[
 (M_M/IM_M)\otimes_{B_B}(B_B[T_a])_k
       \longrightarrow I^kM_M/I^{k+1}M_M
\]
are isomorphisms for every $k$.  Surjectivity is multiplication
by the degree-$k$ monomials, and injectivity is PMS7.
Flatness of $C$ over $B_B$ preserves their exact kernels and
quotients, proving the same statement for $N$.
Also (PMS.13) and right exactness of tensor products give
\[
 N/IN\simeq C\otimes_R X,                        \tag{PMS.23}
\]
as $C$-modules.  Thus the fiber is finite free over the actual
$L_a$ by PMS12.

We prove directly that, if the rank of this fiber is $v$,
\[
 N\simeq L_a[[Y_a:a\in\sigma]]^{,v},\qquad Y_a\mapsto y_a.
                                                        \tag{PMS.24}
\]
First check convergence.  Every finite module over $C$ is complete
and separated for $\mathfrak m_C$.  For every fixed $k$, the
submodule $I^kN$ is closed in that topology: the finite quotient
$N/I^kN$ is separated by Krull's intersection theorem, Stacks
\href{https://stacks.math.columbia.edu/tag/00IP}{00IP}, applied
to $C$ and its proper maximal ideal.  An $I$-adic Cauchy sequence
is $\mathfrak m_C$-adic Cauchy since $I^k\subset\mathfrak m_C^k$;
its limit remains in each prescribed coset of $I^kN$ by that
closedness.  This proves $I$-adic completeness and separation of
$N$ and, by the same reasoning, of $C$ itself.
The formal-series evaluation map $L_a[[Y]]\to C$ is consequently
well defined.  Its coefficients use the fixed map $L_a\to C$.

Choose lifts $v_1,\ldots,v_v$ of a basis of (PMS.23).
For $z\in N$, express its class modulo $IN$ in that basis.
The remaining error lies in $IN$.  Express an error in $I^kN$
as a finite sum $\sum_{|\alpha|=k}y^\alpha z_\alpha$, and reduce
each $z_\alpha$ modulo $IN$ in the same $L_a$-basis.
This supplies degree-$k$ coefficient corrections and leaves an
error in $I^{k+1}N$.  The coefficient sequences converge in
$L_a[[Y]]$, and the module errors converge to zero as just proved.
Hence the map from the right side of (PMS.24) is onto.

For injectivity, suppose a nonzero vector of formal series gives
a relation.  Let $k$ be its least nonzero total degree.  Write
its homogeneous degree-$k$ part as
$\sum_{|\alpha|=k}Y^\alpha(a_{\alpha,1},\ldots,a_{\alpha,v})$.
The remaining terms evaluate into $I^{k+1}N$, so quasiregularity
gives $\sum_i a_{\alpha,i}v_i\in IN$ for each $\alpha$.
Freeness of the fiber over $L_a$ forces every $a_{\alpha,i}=0$,
contrary to the choice of $k$.  This proves injectivity.
If $v=0$, (PMS.23) and local Nakayama give $N=0$ immediately.
Thus every case, including a zero support component, is covered.
\hfill\textit{Source endpoint: \texttt{PmDimQuot.pmDim\_step};
the argument replaces the depth package by a complete basis proof.}
% END PMS13

% BEGIN PMS14
\paragraph{PMS14. Every minimal support component has the required dimension.}
\label{claim:PMS14}
Suppose the original coefficient completion $L_a$ is regular local
of dimension $d$.  Set $A_0=L_a[[Y]]$ and $r=|\sigma|$.
This is a Noetherian regular local domain of dimension $d+r$.
For completeness, choose a regular system of parameters
$z_1,\ldots,z_d$ of $L_a$.
The ring is Noetherian by Stacks
\href{https://stacks.math.columbia.edu/tag/00FM}{00FM}, Lemma
10.31.2.  The successive quotients of $L_a$ by initial parameter
segments are regular local domains by
\href{https://stacks.math.columbia.edu/tag/00NQ}{00NQ} and
\href{https://stacks.math.columbia.edu/tag/00NP}{00NP}.
Consequently the ideals obtained by successively adjoining
$z_1,\ldots,z_d,Y_1,\ldots,Y_r$ form a strict prime chain in
$A_0$ of length $d+r$.  They also generate its maximal ideal.
The generator bound for height, Stacks
\href{https://stacks.math.columbia.edu/tag/0BBZ}{0BBZ}, gives the
reverse inequality; hence $\dim A_0=d+r$ and $A_0$ is regular.
There is no assumption on the characteristic of $L_a$ here.

If $N=0$ there are no primes minimal over its unit annihilator.
Otherwise PMS13 makes $N$ a nonzero finite free $A_0$-module.
Let $C_0=C/\operatorname{Ann}_C N$.
Its faithful action embeds it as an $A_0$-submodule of
$\operatorname{End}_{A_0}N$.  Therefore $C_0$ is finite and
torsion free over $A_0$, and the map $A_0\to C_0$ is injective.
Every minimal prime $Z_0$ of $C_0$ is associated: Stacks
\href{https://stacks.math.columbia.edu/tag/02CE}{02CE} applies
to the finite module $C_0$ over the Noetherian ring $C_0$.
If $0\ne a\in Z_0\cap A_0$, an element whose annihilator is
$Z_0$ would be nonzero $A_0$-torsion in $C_0$, a contradiction.
Thus $Z_0\cap A_0=0$.
The extension $A_0\hookrightarrow C_0/Z_0$ is finite and hence
integral.  Stacks
\href{https://stacks.math.columbia.edu/tag/00OK}{00OK} applies
to this injective integral extension and yields
\[
 \dim(C/Z)=d+r
 \quad\text{for every }Z\text{ minimal over }\operatorname{Ann}_C N.
                                                        \tag{PMS.25}
\]
This is an equality, stronger than the lower bound needed in S5.
It depends on the exact free fiber (PMS.23) and the actual
completed module, not on a dimension assigned to an auxiliary
abstract module.
\hfill\textit{Source: \texttt{PmDimQuot.pmDim}.}
% END PMS14

% BEGIN PMS15
\paragraph{PMS15. Construct every field of the original S5 output.}
\label{claim:PMS15}
Take all hypotheses of S5 in PMS1, with $X$ finite free, an
arbitrary index set $\iota$, primes $q_i$ containing $I$ and
$\mathfrak a$ and avoiding $f$, and, for every $i$, the given
regular coefficient completion of common dimension $d$.
Construct $B_B,M_M$ once, before the points, by PMS5--PMS8.
For each $i$ use $C_i=C_{q_i}$ and $N_i=C_i\otimes_{B_B}M_M$.
Set
\[
 B_i=C_i/\mathfrak aC_i,\qquad
 S_i=(\operatorname{Ann}_{C_i}N_i+\mathfrak aC_i)/\mathfrak aC_i.
                                                        \tag{PMS.26}
\]
These are complete Noetherian local $R/\mathfrak a$-algebras.
Exactness of finite-module completion identifies $B_i$ with the
completion of $(R/\mathfrak a)_{q_i/\mathfrak a}$, and preserves
the structural map.  This is its exact \texttt{cpl} field.

Formula (PMS.16) implies
$\mathfrak aC_i\subset\sqrt{\operatorname{Ann}_{C_i}N_i}$.
Therefore adjoining $\mathfrak aC_i$ to the annihilator changes
neither its containing primes nor its minimal primes.  For a
prime minimal over $S_i$, its inverse image $Z$ in $C_i$ is
minimal over $\operatorname{Ann}_{C_i}N_i$, and
$B_i/(Z/\mathfrak aC_i)=C_i/Z$.
PMS14 gives dimension $d+|\sigma|$ for that quotient.  This
proves the literal \texttt{dim} field.

A prime $Q$ of $B_i$ containing all $y_a$ lifts to a prime of
$C_i$ with the same property.  Formula (PMS.20) gives exactly
\[
 S_i\subset Q
 \quad\Longleftrightarrow\quad
 \operatorname{Ann}_R X
   \subset (Q\cap(R/\mathfrak a))\cap R,          \tag{PMS.27}
\]
which is \texttt{fib}.  Finally lift the two primes $Z\subset Q$
in the original movement antecedent to $C_i$ and apply PMS11.
The resulting $Z'$ in $C_{i'}$ contains $\mathfrak aC_{i'}$ by
(PMS.16), so descends to $B_{i'}$.  Its contraction to
$R/\mathfrak a$ equals that of the original $Z$, because their
contractions to $R$ are equal and both contain $\mathfrak a$.
This proves \texttt{move} with all original quantifiers.

Thus \emph{the complete original S5 is proved}.  It is not
restricted to $h=0$, a finite list of centers, or a minimal-prime
movement assertion.  Its ring/module input hypotheses remain
the precise original hypotheses of PMS1, supplied arithmetically
in the second half of this chapter.
\hfill\textit{Sources: \texttt{OutOfNil},
\texttt{AutomorphicDetection/AnnBc.s5}.}
% END PMS15

% BEGIN PMS16
\paragraph{PMS16. Compute the normalized coefficient dimension directly.}
\label{claim:PMS16}
Let $O$ be the original complete coefficient DVR, $K$ the finite
field of an original point, and suppose $O\to K$ is onto.
Let
\[
 L_0=O[[s,w_1,\ldots,w_D]],\qquad
 g:L_0\longrightarrow K[[t]],\qquad g(s)=t,       \tag{PMS.28}
\]
be the original coefficient map.  It is local: if $b$ were in
the maximal ideal of $L_0$ and $g(b)$ had nonzero constant
coefficient, choose a unit $o\in O$ with that constant image;
$1-o^{-1}b$ would be a unit with nonunit image.  Thus $g$ is
continuous for the maximal-ideal topologies.
The surjection $O\to K$ has kernel $(\varpi)$, where $\varpi$
is a uniformizer.  Indeed a nonzero prime in a DVR is maximal,
and the kernel cannot be zero: an injection of the nonfield
DVR into the finite field $K$ would make it a finite domain and
hence a field.

Write $g(w_i)=f_i(t)\in tK[[t]]$ and lift its coefficients to
$F_i(s)\in sO[[s]]$.
The continuous coordinate change $z_i=w_i-F_i(s)$ is an
automorphism of $L_0$, with inverse $w_i=z_i+F_i(s)$.
Coefficientwise lifting and separation now give
\[
 \ker g=(\varpi,z_1,\ldots,z_D),\qquad
 L_0/\ker g=K[[s]].                              \tag{PMS.29}
\]
There is a strict chain of prime ideals of length $D+1$ ending
at this kernel: start with $(0)$, then adjoin $\varpi$, and
then $z_1,\ldots,z_D$.  The quotients are power-series domains.
These same $D+1$ elements generate the maximal ideal of
$(L_0)_{\ker g}$.  The height bound (tag 0BBZ) gives dimension
exactly $D+1$, and that local ring is regular by the definition
of regular local ring.
Completion preserves dimension and regularity by Stacks
\href{https://stacks.math.columbia.edu/tag/07NV}{07NV} and
\href{https://stacks.math.columbia.edu/tag/07NY}{07NY}.
Consequently the exact coefficient completion in S5 is regular
of dimension $1+D$.

For the source this is $L_0=\mathrm{abelian}_2(O,\mathrm{Unit},D)$;
its single Unit variable is exactly $s$.  The hypothesis that
this variable maps to the original $t$ is indispensable.  For
$\sigma=\mathrm{Fin}(q)\times\mathrm{Fin}(n)$, PMS15 therefore
gives the source dimension
$1+D+nq=1+D_0$.  In the selected Count-II construction this is
$1+4[L^+:\mathbb Q]+4q$.
\hfill\textit{Source: \texttt{LaDimArith.laDim}.}
% END PMS16

\subsection{Supply the hypotheses from the original Taylor--Wiles objects}

% BEGIN PMS17
\paragraph{PMS17. Fix the actual universal ring and all actions.}
\label{claim:PMS17}
The input consists of the original data
\[
 \ell,L,G_0,T,\Lambda,n,O,r,B,u,J,D,A
\]
with $\Lambda:\mathrm{Lim}(T)$ and
$A:\mathrm{TWAut}(u,J,\mathrm{Unit},D,T)$.
Assume $\ell$ prime, $O$ the original coefficient DVR,
$T.\mathrm{ProL}$ and $\mathrm{Adapted}(T,r)$.
For the relation assertions also assume $n\geq1$.
Use the full actual $W:\mathrm{RRObj}(\Lambda,A)$ and
$m:\mathrm{RMaps}(W)$ supplied by PRG.
Write
\[
 \begin{gathered}
 \Phi^0=\ker(\Lambda.G.\eta\circ\Lambda.\pi_o),\quad
 R=W.R,\quad\mathcal D=W.t.D,\\
 S_\infty=L_0[[Y_{ji}]],\quad
 L_0=\mathrm{abelian}_2(O,\mathrm{Unit},D).
 \end{gathered}
                                                        \tag{PMS.30}
\]
The maps $\pi_o^0:\Phi^0\to\Lambda.G^0$ and
$\pi_N^0:\Phi^0\to T.G_N^0$ are the original codrestrictions;
their construction uses the actual common residual identity.
Put $h=\Lambda.h$, $h_N=T.h(N)$, $s=m.s$,
$\rho_N=m.\rho(N)$ and $\pi=m.\pi$.
Let $\mathrm{Sp}=\mathrm{Spo}(A)$ be the original simultaneous
splitting algebra for $u.D$ at $T.sg$, with roots
$u^\mathrm{sp}_{ji}$ reducing to $A.ut_{ji}$.
Let $\mathrm{Sp}_N$ be the corresponding splitting algebra for
$u.D$ at $h_N^0(T.ph_{Nj})$.
Both use the same formal root variables over $B$.

In this notation the actual modules, not substitutes, are
\[
 M_N=(A.A(N)).M,\quad r_0x=\rho_N(r_0)x;
 \qquad X_0=\mathrm{Sp}\otimes_B A.b.M,\quad
 r_0x=\pi(r_0)x.                                  \tag{PMS.31}
\]
The $L_0$-action on $X_0$ is through $A.b.\mathrm{lam}$ and
$B\to\mathrm{Sp}$.  Use the one ultrafilter $\mathscr U=\Lambda.U$.
Its tail condition is a field of this same Lim.
All five components of each original \texttt{Tup.Maps} are
retained: the full determinant law, every local character, the
single section, all roots $U_{ji}$ and all $y_{ji}$.
In particular, for \emph{every} $z\in S_\infty$,
\[
 \rho_N(s(z))=A_N.s(z),\qquad
 \pi(s(z))=\mathrm{base}\bigl(A.b.\mathrm{lam}(z|_{Y=0})\bigr).
                                                        \tag{PMS.32}
\]
These are full power-series identities from RMaps, not identities
on a presumed polynomially dense list alone.

Set $y_{ji}=W.t.y_{ji}$ and define in $R$
\[
 \begin{split}
 f&=\prod_j\prod_i\prod_{i'\ne i}
       (U_{ji}-U_{ji'})(U_{ji}-U_{ji'}),\\
 F_N&=\prod_j\prod_i\prod_{i'\ne i}
       (U_{ji}-U_{ji'})(U_{ji}-T.Nq_{Nj}U_{ji'}).
 \end{split}                                             \tag{PMS.33}
\]
These are exactly $\mathrm{fo}(W)$ and
$\mathrm{fEl}(T.Nq_N,W.t.U)$, including the ordered-pair
indexing.  Let $\mathfrak a=\mathrm{ii}(W)$, generated by the
six original relation families specified in PMS25--PMS27.
No ordinary/control, automorphic realization, or same-$J$
\texttt{HasAut} producer has been assumed proved by this notation.
\hfill\textit{Sources: \texttt{SpecializationMapCompUniversalSection},
\texttt{ThetaCoefficientMapEq}.}
% END PMS17

% BEGIN PMS18
\paragraph{PMS18. The actual special fiber is finite free with the required action.}
\label{claim:PMS18}
The original \texttt{TWBase} says that $M_0=A.b.M$ is finite
free over $L_0$ and its $B$-action agrees with $A.b.\mathrm{lam}$.
PRG's full \texttt{SpoCo} gives a finite free $B$-basis of the
actual $\mathrm{Sp}$.  Tensoring that basis gives an isomorphism
\[
 X_0\simeq\bigoplus_{\text{that basis}} M_0
\]
of $L_0$-modules: on every tensor, the two scalar actions agree
by the original \texttt{TWBase} compatibility.  Thus $X_0$ is
finite free over $L_0$, including the possible rank-zero case.
Formula (PMS.32) with $z\in L_0$ identifies this coefficient
action with $c=s|_{L_0}$; with $z=Y_{ji}$ it gives $y_{ji}X_0=0$.
These verify every coefficient and fiber hypothesis in PMS1.

The ring $L_0$ is complete Noetherian local with finite residue
field: it is a finite-variable formal power-series ring over
the original complete coefficient DVR $O$.  Its maximal ideal
is generated by the uniformizer and all its variables; its
finite residue field is the one supplied by $\mathrm{IsCoef}$.
The uniformizer-adic completeness of $O$, together with formal
coefficient expansions, gives maximal-ideal completeness of
$L_0$.  No assertion that $B$ itself is free over $L_0$ is used.
\hfill\textit{Source: \texttt{PatchedModuleFreeFiniteOverIwasawaAlgebra}.}
% END PMS18

% BEGIN PMS19
\paragraph{PMS19. Every required finite quotient is free at all sufficiently large levels.}
\label{claim:PMS19}
The original \texttt{TWRing} contains integers $a_{Nj}\geq N$
and an additive isomorphism
\[
 M_N\simeq
 \bigl(S_\infty/\Delta_N\bigr)^{r_N},\qquad
 \Delta_N=\bigl((1+Y_{ji})^{\ell^{a_{Nj}}}-1\bigr)_{j,i},
                                                        \tag{PMS.34}
\]
intertwining the entire $S_\infty$-action.  Fix $k,j_0\geq0$.
For $N\geq j_0+k$ we have
\[
 \Delta_N\subset (Y)^k+\mathfrak m_{L_0}^{j_0}S_\infty.
                                                        \tag{PMS.35}
\]
If $k=0$ or $j_0=0$ the right side is the unit ideal.
Otherwise expand each binomial.  Terms of degree at least $k$
belong to $(Y)^k$.  For $1\leq v<k$ the identity
\[
 v\binom{\ell^a}{v}=\ell^a\binom{\ell^a-1}{v-1}
\]
shows that the integer coefficient is divisible by
$\ell^{a-v_\ell(v)}$, hence by $\ell^{j_0}$, since
$a\geq j_0+k$ and $v_\ell(v)\leq v-1<k$.
The coefficient is zero if $v>\ell^a$; that case also satisfies
the assertion.  As $\ell\in\mathfrak m_{L_0}$, this proves
(PMS.35) in the original coefficient ring without dividing by
$v$ in that ring.

Quotient (PMS.34) by the image of the right-side ideal.
Using (PMS.32) to identify its action with the $R$-action gives
exactly (PMS.2), with the same $r_N$ and the same $M_N$.
Thus the assertion holds on the tail $N\geq j_0+k$, and hence
on a $\Lambda.U$-large set, for every $k,j_0$.
No bound on $r_N$ is asserted or required.
\hfill\textit{Source: \texttt{FreeNQuot.freeNQuot}.}
% END PMS19

% BEGIN PMS20
\paragraph{PMS20. Full determinant comparison and the universal-ring uniqueness test.}
\label{claim:PMS20}
We use the integral word identity proved in JSP9: each
characteristic coefficient of a finite sum of group elements,
with arbitrary scalar coefficients, is an integral polynomial
in characteristic coefficients of group words.  The identity
is natural under every commutative scalar extension.  In
particular equality of all $\lambda_i([g])$ implies equality
of the full determinant laws.  This statement uses no division
by $n!$ and is valid for $n=4,\ell=3$.

Let $\mathcal E$ be a continuous determinant law on a profinite
group $H$, valued in a ring with the specified ideal-adic
coefficient topology.  For each power $\mathfrak b^j$ there is
one open normal $V\subset H$ such that
\[
 \overline{x}=\overline{x'}\text{ in }R[H/V]
 \ \Longrightarrow\
 \lambda_i^{\mathcal E}(x)-\lambda_i^{\mathcal E}(x')
      \in\mathfrak b^j\quad\hbox{for all }i,x,x'.             \tag{PMS.36}
\]
Indeed each of the finitely many functions
$g\mapsto\lambda_i([g])\bmod\mathfrak b^j$, $0\leq i\leq n$,
is continuous to a discrete space.  Compactness gives a finite
clopen fiber partition.  A common open normal subgroup has
cosets refining all these partitions: choose sufficiently small
translated neighborhoods on a finite cover, and then take a
normal subgroup contained in the finitely many corresponding
identity neighborhoods.  The functions are therefore constant
on its cosets.  Replacing every group element in a finite sum
by its coset representative leaves all characteristic
coefficients unchanged modulo $\mathfrak b^j$, by the integral
word identity; combining equal cosets proves (PMS.36).
The source quantifies over every natural $i$, including $i>n$:
its definition uses the truncated subtraction $n-i=0$, so
$\lambda_i=(-1)^{i-n}\lambda_n$ in that range.  These signed
constant-coefficient values are already covered by the same
subgroup; no extra continuity or vanishing claim is needed.
The identity for all coefficient algebras follows from the same
integral polynomials.  This proves the full \texttt{dSees}
scope.  Applied to $H=\Phi^0$, refine $V$ to an open normal
subgroup of the ambient $\Lambda.\Phi$ contained in $\Phi^0$;
the normal core is a finite intersection because $\Phi^0$ is
open.  The property persists on this smaller subgroup.  This
is the exact \texttt{phSees} used below.

There is also the following exact \texttt{rrExt} consequence.
Two $O$-algebra maps $b,b':R\to C$, where $C$ is an original
$\mathrm{IsCO}$ target, are equal if they agree on all
$\mathcal D.\lambda_i([g])$, all $W.t.\psi_{vi}(g)$, the
section $W.t.s$, every $U_{ji}$ and every $y_{ji}$.
The integral identity gives equality of the full determinant
component.  The other equalities give the remaining four tuple
components.  The push-forward tuple is a genuine \texttt{Tup}:
the original local, continuous $O$-algebra maps preserve the
residual law, character restrictions and maximal-ideal
conditions, as in PRG.  Both maps therefore satisfy the same
unique-map clause of this actual $W.\mathrm{univ}$, and are
equal.  This uses its complete universal property, not an
unproved generation statement about traces alone.
\hfill\textit{Sources: \texttt{LamEqOfKer.dSees}, \texttt{PhSees},
\texttt{RrExt.rrExt}; inherited input JSP9.}
% END PMS20

% BEGIN PMS21
\paragraph{PMS21. One event compares the specialization maps on every ring element.}
\label{claim:PMS21}
For each $v\geq0$, there is one $\Lambda.U$-large event on
which both of the following hold.  Let $P_j$ be the
characteristic polynomial of $u.D$ at $T.sg_j$, and $P_{Nj}$
the one at $h_N^0(T.ph_{Nj})$.  Then every coefficient of
$P_j-P_{Nj}$ belongs to $\mathfrak m_B^v$.  Moreover, for
\emph{every} $a\in R$ and every pair of formal root-series
lifts $z,z'\in B[[Z_{ji}]]$ of $\pi(a)$ and
$\theta_N\rho_N(a)$, respectively,
\[
 z-z'\in\mathrm{spIdeal}(P)+\mathfrak m_B^vB[[Z]].             \tag{PMS.37}
\]
Here $\theta_N$ is the original map in $A.m(N)$.

Choose an ambient open normal subgroup $V\subset G_0.\Gamma$
using (PMS.36) for $u.D$ at precision $v$.
The source $T.ph\_lim$ applies to this $V$: its residual-kernel
guard is satisfied because $T.\mathrm{ProL}$ makes that
relative pro-$\ell$ residual zero.  Intersect the finitely many
events for $j$ with the one uniform event $\Lambda.p\_lim(V)$.
The first assertion follows.  Set
\[
 C_v=B[[Z]]/
      \bigl(\mathrm{spIdeal}(P)+\mathfrak m_B^vB[[Z]]\bigr).
\]
The congruent coefficients make both $\mathrm{Sp}$ and
$\mathrm{Sp}_N$ map to this same quotient, sending each root
to the same $A.ut_{ji}+Z_{ji}$.  If the quotient is the zero
ring, (PMS.37) is automatic.  Otherwise PRG's \texttt{SpoCo}
and the proper-quotient completion facts make it an
$\mathrm{IsCO}(O)$ target.

Compare the composites $R\to\mathrm{Sp}\to C_v$ and
$R\xrightarrow{\rho_N}R_N\xrightarrow{\theta_N}\mathrm{Sp}_N
\to C_v$ by PMS20.  Their group characteristic coefficients
agree for \emph{all} $g\in\Phi^0$ on the one uniform
$\Lambda.p\_lim(V)$ event.  Their characters agree by the
literal $\theta_N.\mathrm{pair}$ field.  Their section and
$y$ components agree by both all-series identities in
(PMS.32) and the original $\theta_N$ coefficient identity;
their root components agree by $\theta_N(U_{Nj i})=u^{\rm sp,N}_{ji}$.
Thus the two entire ring maps are equal on that event.
Applying their equality to any $a$ and any lifts proves
(PMS.37).  The order is
$\forall v\ \forall^{\mathscr U}N\ \forall a,z,z'$;
there is no point-dependent large event hidden in the argument.
\hfill\textit{Source: \texttt{PiCong.piCong}.}
% END PMS21

% BEGIN PMS22
\paragraph{PMS22. Compare splitting-module quotients and the exact scalar factor.}
\label{claim:PMS22}
Suppose $P_j\equiv P'_j\pmod{H}$ over $B$, and $H$ kills
$M_0/AM_0$, where $A$ is a coefficient ideal acting on $M_0$.
Both modules
\[
 (\mathrm{Sp}(P)\otimes_B M_0)/A,\qquad
 (\mathrm{Sp}(P')\otimes_B M_0)/A
\]
are quotients of the same $B[[Z]]\otimes_B(M_0/AM_0)$.
Their defining root relations differ by elements of
$HB[[Z]]$, which act as zero on that module.  The two quotient
submodules are therefore equal, giving a canonical additive
isomorphism $\kappa$ in either direction.
It intertwines multiplication by every common series in
$B[[Z]]$.  Also $\mathrm{spIdeal}(P)+HB[[Z]]$ kills the first
quotient.  This proves precisely the comparison required from
\texttt{SpCong}; it does not require flatness of $M_0$ over $B$.

For the scalar factor in (PMS.33), if $N\geq j$ then
$T.Nq_{Na}\equiv1\pmod{\ell^j}$ for every $a$.
Reducing the defining finite product modulo $(\ell^j)$ gives
\[
 f-F_N=\ell^j c_N\quad(c_N\in R),\qquad
 \rho_N(F_N)=\mathrm{fEl}(T.Nq_N,A_N.U).          \tag{PMS.38}
\]
In particular the difference lies in $\mathfrak m_R^j$.
The stronger displayed divisibility is needed for the module
quotients: $\ell^j$ comes from
$\mathfrak m_{L_0}^j$, so the difference acts as zero modulo
that coefficient ideal as well.  Merely knowing membership
in $\mathfrak m_R^j$ would not justify that step.
\hfill\textit{Sources: \texttt{SpCong}, \texttt{FoCong.foCong}.}
% END PMS22

% BEGIN PMS23
\paragraph{PMS23. Construct the required paired maps on the original finite quotients.}
\label{claim:PMS23}
Fix $j\geq0$.  The quotient
$M_0/\mathfrak m_{L_0}^jM_0$ is a finite set by PMS18.
The finite-module Nakayama argument in PMS8, now for the
original local ring $B$, gives an integer $v$ with
\[
 \mathfrak m_B^vM_0\subset\mathfrak m_{L_0}^jM_0.             \tag{PMS.39}
\]
Intersect the event of PMS21 at this $v$ with $N\geq j$.
PMS22 produces the canonical map
\[
 \kappa_N:X_0/\mathfrak m_{L_0}^jX_0
   \simeq(\mathrm{Sp}_N\otimes_B M_0)/
                        \mathfrak m_{L_0}^j.
\]
It is $R$-linear for the actions through $\pi$ and
$\theta_N\rho_N$.  To verify this for an arbitrary $a\in R$,
lift its two images to series $z,z'$ as in PMS21.
Their difference belongs to the common ideal in (PMS.37),
which kills this quotient by (PMS.39) and PMS22.
Thus multiplication by the two images agrees.  This checks
every $R$-scalar on the one event.

The original \texttt{TWMaps} supplies additive maps
\[
 \varphi_N:\mathrm{Sp}_N\otimes_B M_0\longrightarrow M_N/IM_N,
 \qquad
 \psi_N:M_N/IM_N\longrightarrow\mathrm{Sp}_N\otimes_B M_0,
\]
semilinear for $R_N\xrightarrow{\theta_N}\mathrm{Sp}_N$,
with composites given by the original finite-level factor in
(PMS.38) and its image under $\theta_N$.
Formula (PMS.32) and the $\theta_N$ coefficient identity imply
that these maps preserve the coefficient submodules
$\mathfrak m_{L_0}^j$; hence they descend to the displayed
quotients.  Compose them with $\kappa_N$ and its inverse.
They are $R$-linear and give maps in both directions between
\[
 X_0/K_{1,j}X_0
 \quad\hbox{and}\quad M_N/K_{1,j}M_N.
\]
Their composites are $f$: the original composite is $F_N$
under the relevant action, and (PMS.38) makes $f-F_N$ act
as zero on these quotients.  This proves every field of the
second finite-level antecedent in S5, with
$\forall j\ \forall^{\Lambda.U}N\ \exists\varphi,\psi$.
No compatible choices of these maps across $j$ or $N$ are
needed, by PMS3.
\hfill\textit{Source: \texttt{EventuallyExistsPatchedModuleQuotientMaps}.}
% END PMS23

% BEGIN PMS24
\paragraph{PMS24. Lift close elements while preserving their original cosets.}
\label{claim:PMS24}
Let $\alpha:\Phi\to H_N$ and $\beta:\Phi\to H_\infty$
be the actual surjections, and let $V\triangleleft\Phi$
satisfy $\ker\alpha\,V=\ker\beta\,V$.
If $x\in\ker\beta$, then $x=kv$ for
$k\in\ker\alpha$, $v\in V$; thus $x^{-1}k=v^{-1}\in V$.
This gives the kernel part of \texttt{closeLift}.

For the other part, the source
$\mathrm{Close}(\alpha,\beta,V,a,b)$ means that there are
$g_0,g'_0\in\Phi$ with
$\alpha(g_0)=a$, $\beta(g'_0)=b$ and $g_0^{-1}g'_0\in V$.
Given a specified lift $g$ of $b$, write
$g_0'^{-1}g=kv$ as above, and put $g'=g_0k$.
Then $\alpha(g')=a$ and
\[
 g^{-1}g'=v^{-1}k^{-1}(g_0'^{-1}g_0)k\in V,
\]
by normality.  Thus the new lift is close to the specified old
lift, not just to some untracked lift.
Apply this with $\alpha=\Lambda.\pi_N$,
$\beta=\Lambda.\pi_o$ and the one subgroup from PMS20.
The required kernel equality is the actual
$\Lambda.\mathrm{ker\_lim}(V)$ event.  Since $V\subset\Phi^0$,
the new lift of an element of $\Phi^0$ remains in $\Phi^0$.
The simultaneous finite lists of Frobenius, tame and local
elements below use finite intersections of the original events.
\hfill\textit{Source: \texttt{CloseLift.closeLift}.}
% END PMS24

% BEGIN PMS25
\paragraph{PMS25. Approximate the kernel and polarization relations.}
\label{claim:PMS25}
Fix a relation generator $a\in R$ and a precision $j$.
We first treat the source constructors \texttt{ker} and \texttt{pol}.
For
\[
 a=\operatorname{tr}_{\mathcal D}([h]([k]-1)),
 \qquad \pi_o^0(k)=1,
\]
choose the subgroup $V$ of PMS20 at precision $j$ and use the
kernel part of PMS24 on the $\Lambda.\mathrm{ker\_lim}(V)$
event.  It gives $k'\in\Phi^0$ with
$\pi_N^0(k')=1$ and $k^{-1}k'\in V$.
Replace $k$ by $k'$ to define $a'$.
The two group-algebra expressions have the same image modulo
$V$, so $a-a'\in\mathfrak m_R^j$ by (PMS.36).
The full determinant component of $\rho_N.\mathrm{Maps}$
gives $\rho_N(a')=0$, since $[k']-1$ maps to zero.
The trace of zero is zero for $n\geq1$, directly from its
characteristic polynomial.

For \texttt{pol}, the original expression is
\[
 a=\lambda_i^{\mathcal D}
       \bigl([g'^{-1}]\epsilon_\infty(g)^{1-n}\bigr)
       -\lambda_i^{\mathcal D}([g]),\qquad
 g'=\Lambda.ct\,g\,\Lambda.ct^{-1}.              \tag{PMS.40}
\]
Here $\epsilon_\infty(g)$ is the unit
$\Lambda.G.\mathrm{frame.eps}(O,R,\pi_o^0g)$.
Let $\epsilon_N(g)$ be the corresponding original cyclotomic
unit from $(T.G_N)^0$ evaluated at $\pi_N^0g$, still in $R$.
The kernel of the original global cyclotomic character modulo
$\mathfrak m_R^j$ is open: the coefficient map kills
$\ell^j$, so it factors through the finite cyclotomic quotient
modulo $\ell^j$ (the zero quotient case is harmless).
The uniform outgoing convergence $\Lambda.p\_lim$ applied to
that kernel gives
$\epsilon_\infty(g)\equiv\epsilon_N(g)\pmod{\mathfrak m_R^j}$.
Their integer powers, including the negative exponent $1-n$,
agree in the unit group of the quotient.

Define $a'$ by replacing only $\epsilon_\infty(g)^{1-n}$ in
(PMS.40) by $\epsilon_N(g)^{1-n}$.
Naturality of the full determinant law makes
$a-a'\in\mathfrak m_R^j$.
The exact identities $\Lambda.\pi_N(\Lambda.ct)=T.G_N.c$
and $g'=\Lambda.ct\,g\,\Lambda.ct^{-1}$ turn
$\rho_N(a')$ into the original polarization identity
$A_N.p.P$, so $\rho_N(a')=0$.
There is no approximation of the conjugating element or its
exponent, and no choice of a new polarization.
\hfill\textit{Sources: \texttt{PatchedTraceSingleMulSubOneApproxAnnihilated},
\texttt{PatchedLamConjTwistApproxAnnihilated}.}
% END PMS25

% BEGIN PMS26
\paragraph{PMS26. Approximate every ordered ordinary and inertia relation.}
\label{claim:PMS26}
For \texttt{ord} fix an original place $v$, a tuple
$g_i\in L.G_v$, lifts $b_i\in\Phi^0$ with
$\pi_o^0(b_i)=\Lambda.G.\phi_v(g_i)$, and $h\in\Phi^0$.
Its generator is
\[
 a=\operatorname{tr}_{\mathcal D}
 \left([h]\prod_{i=0}^{n-1}
             ([b_i]-\psi_{vi}(g_i))\right).       \tag{PMS.41}
\]
The product has the original order of \texttt{List.ofFn}.
Choose $V$ by PMS20.  On the intersection of
$\Lambda.\mathrm{ker\_lim}(V)$ and the uniform
$\Lambda.\phi\_lim(V)$ events, PMS24 supplies
$b_i'\in\Phi^0$ in the same $V$-coset as $b_i$ with
$\pi_N^0(b_i')=(T.G_N).\phi_v(g_i)$.
Replace just these group elements in (PMS.41), leaving all
the universal characters and the order unchanged.
Equation (PMS.36) proves congruence at precision $j$.
Under $\rho_N$ the full character and determinant identities
identify $a'$ with $A_N.p.\mathrm{Ord}$ on the same tuple and
the element $\pi_N^0h$.  Hence $\rho_N(a')=0$.
Nothing commutes the factors in this noncommutative group algebra.

For \texttt{ine} fix $w$, $g\in L.J_u(w)$, and a lift $b$
of $\Lambda.G.\phi_u(w,g)$.  The generator is
$a=\operatorname{tr}_{\mathcal D}([h]([b]-1))$.
Use $\Lambda.\phi_u\_lim(V)$ in place of $\phi\_lim$ to get
a same-coset lift $b'$ of $(T.G_N).\phi_u(w,g)$.
The same calculation gives $a\equiv a'\pmod{\mathfrak m_R^j}$.
The new local image is in the exact residual inertia subgroup
$A_N.r_N.Ip(w)$, not merely the geometric inertia image.
Indeed it is in $(T.G_N).I_u(w)$ by $g\in L.J_u(w)$.
Every residual character on it equals the original residual
character on $G_0.\phi_u(w,g)$ by $A_N.\mathrm{compat}$.
Adaptedness says $r.Ip(w)=G_0.I_u(w)$, so all these values
are one.  This verifies both parts of the defining intersection
$A_N.r_N.Ip(w)=(T.G_N).I_u(w)\cap\bigcap_i\ker\bar\chi_{N,i}$.
The original field $A_N.p.I$ now gives $\rho_N(a')=0$.
These arguments prove the assertions for every original local
place and tuple, at every precision.
\hfill\textit{Sources: \texttt{PatchedTraceLocalCharPolyApproxAnnihilated},
\texttt{EventuallyTraceInertiaCongruenceOfAdapted}.}
% END PMS26

% BEGIN PMS27
\paragraph{PMS27. Approximate the two auxiliary-prime relation families.}
\label{claim:PMS27}
For an auxiliary index $a$, put
\[
 q_1(p,U_a)=\prod_{i=0}^{n-1}([p]-U_{ai}),\qquad
 q_2(p,t,U_a,i,y_{ai})
   =\left(\prod_{\substack{k=0\\k\ne i}}^{n-1}([p]-U_{ak})\right)
                         ([t]-1-y_{ai}).          \tag{PMS.42}
\]
The indicated products retain the exact source order.
For the \texttt{q1} generator
$a_1=\operatorname{tr}_{\mathcal D}([h]q_1(p,U_a))$,
where $\pi_o^0p=\Lambda.ph_a$, use
$\Lambda.\mathrm{ker\_lim}(V)$ and
$\Lambda.ph\_lim(V)$ with PMS24 to obtain a lift $p'$ in
the same $V$-coset with $\pi_N^0p'=T.ph_{Na}$.
Let $a_1'$ be the expression with $p'$.
PMS20 gives $a_1-a_1'\in\mathfrak m_R^j$.
The five exact tuple identities then give
\[
 \rho_N(a_1')
  =\operatorname{tr}_{A_N.p.D}
       ([\pi_N^0h]q_1(T.ph_{Na},A_N.U_a)).
\]
The original $\mathrm{TWMaps.rel1}$ says that
$(\rho_N(F_N)\rho_N(a_1'))^{A.ex}$ kills the actual $M_N$.

For \texttt{q2}, there is also the original lift $t$ with
$\pi_o^0t=\Lambda.ta_a$.
Intersect in addition with $\Lambda.ta\_lim(V)$ and lift both
$p,t$ in their specified cosets to $p',t'$ with
$\pi_N^0p'=T.ph_{Na}$ and $\pi_N^0t'=T.ta_{Na}$.
The full expressions modulo $V$ agree, so the corresponding
$a_2'$ is congruent to $a_2$.
Under $\rho_N$ it is the exact finite-level ordered $q_2$
expression: both $U$ and $y$ are carried by the original tuple.
The original $\mathrm{TWMaps.rel2}$ therefore gives
$(\rho_N(F_N)\rho_N(a_2'))^{A.ex}M_N=0$.

The exponent is the same $A.ex$ stored once in TWAut, for every
level, auxiliary index, $i$, and test element $h$.
There is no new level-dependent exponent and no use of a
generic semisimple quotient in place of the given module.
\hfill\textit{Sources: \texttt{MapRingHomSingleMulCharpolyElement.t3Q1},
\texttt{MapRingHomSingleMulPartialCharpolyElement.t3Q2}.}
% END PMS27

% BEGIN PMS28
\paragraph{PMS28. One exponent supplies the whole approximate-annihilator ideal.}
\label{claim:PMS28}
PMS25--PMS27 give, for every one of the six literal generators
$a$ and every $j$, on a $\Lambda.U$-large event an $a'$ with
\[
 a-a'\in\mathfrak m_R^j,\qquad
 \rho_N(a')=0
 \quad\hbox{or}\quad
 (\rho_N(F_N)\rho_N(a'))^{A.ex}M_N=0.             \tag{PMS.43}
\]
Intersect with the tail $N\geq j$ and set
\[
 b=(F_Na')^{A.ex+1}\in R.
\]
The finite product identity (PMS.38) shows
$fa-F_Na'\in\mathfrak m_R^j$.
The difference of equal positive powers is divisible by the
difference of their bases, so
\[
 (fa)^{A.ex+1}-b\in\mathfrak m_R^j.              \tag{PMS.44}
\]
On $M_N$ multiplication by $b$ is its image under $\rho_N$.
In the first case of (PMS.43) it is zero because $A.ex+1>0$.
In the second it is one extra multiple of the annihilating
$A.ex$ power, so is also zero.  This includes $A.ex=0$.

Thus every generator satisfies the literal
$\mathrm{IsApproximateAnnihilatorAtLevels}(m,a)$, with the
single exponent $e=A.ex+1$ chosen before the generator,
precision or level.  Take $G_s=\{a:\mathrm{Rel}(W,a)\}$;
by definition its ideal is $\mathrm{ii}(W)$.
This supplies the complete third antecedent (PMS.3).
The large event may depend on a relation and precision,
exactly as allowed; no intersection over infinitely many
relations was taken.
\hfill\textit{Source: \texttt{T3Of.t3Of} and all six constructors
of \texttt{Rel}.}
% END PMS28

% BEGIN PMS29
\paragraph{PMS29. Install the actual patched module and the full support output.}
\label{claim:PMS29}
For the fixed original data of PMS17, PMS18, PMS19, PMS23 and
PMS28 verify all inputs of PMS1 with
\[
 \begin{gathered}
 (L_0,R)=(\mathrm{abelian}_2(O,\mathrm{Unit},D),W.R),\\
 (\sigma,\mathscr U)
   =(\mathrm{Fin}(T.q)\times\mathrm{Fin}(n),\Lambda.U),\\
 (f,\mathfrak a,X,M_N)
   =(\mathrm{fo}(W),\mathrm{ii}(W),X_0,\mathrm{MN}(m,N)).
 \end{gathered}
                                                        \tag{PMS.45}
\]
Consequently PMS5--PMS8 produce an original
$\mathrm{PMod}(W.t.y,\mathrm{fo}(W),\mathrm{ii}(W),X_0)$.
It is constructed from these very $M_N$ with their actions
through $\rho_N$, and from this one $\Lambda.U$.
In particular its fiber is $\mathrm{Spo}(A)\otimes_B A.b.M$
with its action through the same $\pi$.

For any family of primes $q_i\subset W.R$ containing all
$y_{ji}$ and $\mathrm{ii}(W)$ and avoiding $\mathrm{fo}(W)$,
if the coefficient completions have common regular dimension
$d$, PMS15 constructs
\[
 \mathrm{PatchedSupportCompletions}
   \bigl(s,\mathrm{ii}(W),\operatorname{Ann}_{W.R}X_0,
          q,d+nT.q\bigr).                        \tag{PMS.46}
\]
Every field, including movement between arbitrary centers,
is supplied.  For the normalized original points whose
coefficient maps satisfy (PMS.28), PMS16 makes
$d=1+D$ and the dimension parameter is exactly
$1+D+nT.q=1+D_0$.
No new coefficient extension, ideal, level system, point or
ultrafilter has been selected in this step.

The possible zero module causes no exception: its support
ideal is the unit ideal and the dimension statement is
vacuous; the fiber equivalence and movement implications
remain correct.  A later starting support prime, when
provided by the original argument, automatically excludes
that case at its center.
\hfill\textit{Source: the actual S5 application in
\texttt{AutomorphicDetection/LiftPt.core}.}
% END PMS29

\subsection{The exact maps to the original Hecke ideal}

% BEGIN PMS30
\paragraph{PMS30. An integral trace identity proves the full kernel inclusion.}
\label{claim:PMS30}
We first prove the determinant identity needed for \texttt{q1},
without assuming a matrix representation or inverting $n!$.
Let $D$ be a homogeneous multiplicative determinant law of
positive degree $n$ on an algebra over a commutative ring $A$.
In the dual numbers,
$D(1+\varepsilon z)=1+\varepsilon\operatorname{tr}_D(z)$.
Multiplicativity applied to
$(1+\varepsilon z)(1+\varepsilon w)$, together with scalar
naturality, proves that $\operatorname{tr}_D$ is $A$-linear.
For any $x,h$, work in
$A[t,s]/(t^{n+1},s^2)$; all inverses in the next expression
are finite geometric sums.  The equality
\[
 1-tx-sh=\bigl(1-sh(1-tx)^{-1}\bigr)(1-tx)
\]
and multiplicativity give
\[
 D(1-tx-sh)
   =\left(1-s\sum_{k=0}^{n}t^k\operatorname{tr}_D(hx^k)\right)
       \left(\sum_{i=0}^{n}(-1)^i\lambda_i^D(x)t^i\right).
\]
The coefficient of $st^n$ on the left is zero: a degree-$n$
polynomial law evaluated on $1-tx-sh$ has total degree at
most $n$ in $t,s$.
On the right that coefficient is
$-\operatorname{tr}_D(h\chi_x^D(x))$.
Therefore, over every scalar algebra,
\[
             \operatorname{tr}_D(h\chi_x^D(x))=0.             \tag{PMS.47}
\]
This calculation is integral and applies at $\ell=3,n=4$.

We now prove the literal inclusion
\[
                  \mathrm{ii}(W)\subset\ker\pi.              \tag{PMS.48}
\]
Push group-algebra expressions using the full determinant
component of $\pi.\mathrm{Maps}$ and the group map
$h^0\pi_o^0$.
For \texttt{ker}, $[k]-1$ becomes zero.
For \texttt{pol}, the same conjugation and cyclotomic
compatibilities identify the image with $u.P$, including
the inverse and exponent $1-n$.
For \texttt{ord}, the character component and local
compatibility turn the exact ordered product into $u.\mathrm{Ord}$.
For \texttt{ine}, adaptedness puts the local image of $L.J_u$
inside $r.Ip$, so $u.I$ gives zero.
For \texttt{q1}, $h^0(\Lambda.ph_j)=T.sg_j$ and the roots
in $\mathrm{Sp}$ split precisely the characteristic polynomial
of $u.D$ at that element.  Thus the image of $q_1$ is
$\chi_x(x)$ for $x=[T.sg_j]$, and (PMS.47) kills its trace
after every left multiplier.
For \texttt{q2}, the final factor maps to
$[h^0(\Lambda.ta_j)]-1-\pi(y_{ji})=1-1-0=0$.
All six generators therefore map to zero, and so does their
ideal.  This proves \texttt{piKer} in its complete original
scope $n\geq1$ and adapted data; no ProL assumption is needed
for this particular inclusion.
\hfill\textit{Source: \texttt{CharpolyElementEqAevalGroupElement.piKer}.}
% END PMS30

% BEGIN PMS31
\paragraph{PMS31. The specialization onto the original splitting algebra is surjective.}
\label{claim:PMS31}
Assume the exact source generation hypothesis
\[
 \mathrm{GenBy}\ O\
 \left(\mathrm{pairVals}(u)\cup
       \{A.b.\mathrm{lam}(X_{\mathrm{Unit}})\}\right).
                                                        \tag{PMS.49}
\]
Every displayed generator of $B$ has its image in
$\pi(R)\subset\mathrm{Sp}$.
For a group characteristic coefficient use surjectivity of
$h^0\pi_o^0$ and the full determinant component of
$\pi.\mathrm{Maps}$.
For a character value use its character component.
For the original Unit section use its section component.
Every root $u^\mathrm{sp}_{ji}$ is $\pi(U_{ji})$.
The full \texttt{SpoCo} statement from PRG says that these
roots generate $\mathrm{Sp}$ as a $B$-algebra.

The algebraic subring $\pi(R)$ is dense in $\mathrm{Sp}$ for
its maximal-ideal topology.
Indeed (PMS.49) approximates any $b\in B$ by polynomials over
$O$ in the displayed generators, at any prescribed precision.
The local map $B\to\mathrm{Sp}$ takes
$\mathfrak m_B^k$ into $\mathfrak m_{\mathrm{Sp}}^k$.
Polynomials in the roots have the same approximation property,
and they are all of $\mathrm{Sp}$ by its algebraic generation.

We prove that density here gives actual surjectivity.
The map $\pi$ is a local $O$-algebra map between complete
Noetherian local rings with the same residue field as $O$.
Choose a basis of
$\mathfrak m_{\mathrm{Sp}}/\mathfrak m_{\mathrm{Sp}}^2$
and approximate its finitely many lifts by $\pi(a_i)$ modulo
$\mathfrak m_{\mathrm{Sp}}^2$.
The residue identity forces $a_i\in\mathfrak m_R$.
By Nakayama their images generate
$\mathfrak m_{\mathrm{Sp}}$.
For any $z\in\mathrm{Sp}$, choose a coefficient in $O$ with
its residue.  Successively express the degree-$k$ error
modulo $\mathfrak m_{\mathrm{Sp}}^{k+1}$ as a homogeneous
degree-$k$ polynomial in the $\pi(a_i)$ with residue
coefficients lifted to $O$.
Lift this polynomial to $R$; it lies in $\mathfrak m_R^k$.
The correction series converges in $R$, and its image is
$z$, since a local map preserves these powers and the target
is separated.  Thus
\[
                         \pi:R\twoheadrightarrow\mathrm{Sp}. \tag{PMS.50}
\]
This is the full \texttt{piSurj}, not a surjection after
inverting $f$ or after reducing modulo the maximal ideal.
\hfill\textit{Source: \texttt{PiSurj.piSurj}.}
% END PMS31

% BEGIN PMS32
\paragraph{PMS32. The support condition is exactly the original same-ideal condition.}
\label{claim:PMS32}
For every prime $P\subset\mathrm{Sp}$, the conclusions of
PMS31 and the fixed action in (PMS.31) give
\[
 \begin{split}
 \operatorname{Ann}_{R}X_0\subset\pi^{-1}P
 &\ \Longleftrightarrow\
 \operatorname{Ann}_{\mathrm{Sp}}X_0\subset P\\
 &\ \Longleftrightarrow\
 \operatorname{Ann}_{B}M_0\subset P\cap B\\
 &\ \Longleftrightarrow\ J\subset P\cap B .
 \end{split}                                             \tag{PMS.51}
\]
The first step uses the literal identity
$\operatorname{Ann}_R X_0=\pi^{-1}(\operatorname{Ann}_{\mathrm{Sp}}X_0)$
and surjectivity of $\pi$; surjectivity is why no converse is
lost here.  The second is the base-change support calculation
(PMS.19) for the finite $B$-module $M_0$.  The third is the
original $\mathrm{TWBase.supp}$ field at the prime $P\cap B$.
It is the very ideal $J$ in the fixed $A:\mathrm{TWAut}(u,J,\ldots)$.
No radical replacement of the output ideal, unrelated Hecke
ideal, or separately selected automorphic module appears.

Here is the exact completed-support interface for the remaining
one-point argument.  Given an original lifted point
$\xi:\mathrm{Sp}\to K[[t]]$ extending the original $B$-point,
with pairwise distinct root images and the original Unit
section sent to $t$, put $q=\ker(\xi\pi)$.
PMS30 makes $\mathrm{ii}(W)\subset q$.
Every $y_{ji}$ belongs to $q$ by $\pi.\mathrm{Maps}$.
The root distinctness and (PMS.33) give $f\notin q$:
its image is a finite product of nonzero differences in the
domain $K[[t]]$.
If the original map $O\to K$ is onto, PMS16 and PMS29 give
the actual completed support ring with parameter $1+D_0$.
Formula (PMS.51) converts its fiber condition to the original
$J$ at every relevant prime of $\mathrm{Sp}$.
These are precisely the support inputs of CRW18--CRW20.
The existence and every-field matching of the lifted
\texttt{OoPair}/\texttt{OoPt}, and the final selected Core
assembly, are separate PRG consumer arguments; they have
not been inferred solely from the existence of S5.

For Count II the selected CCH--CFO construction has
$h=0$, $a=n=4$, $k.g=0$ and
$D_0=4[L^+:\mathbb Q]+4T.q$.
It applies at every permitted odd $\ell$, including $\ell=3$
whenever the original input hypotheses hold.
All of the constructions above apply on its one actual
$T$, Lim, TWAut, point, ideal and module.
They make no assertion that $\mathrm{Sel1}=0$ or that
arbitrary Kap data satisfy the selected bounds.
The fixed-$\iota$ outer coefficient threshold must still
work for both original rank-four Links and must produce
Hida properties and automorphic support on the same ideal
at that threshold.  The unchanged ordinary/control,
type/neatness and same-$J$ automorphic suppliers, with
the original CohP and LocCFT hypotheses, remain explicit.
Full S5 and the arithmetic construction on a supplied
TWAut remove the patched-module/support producer from
that residual list; they do not by themselves close
\texttt{COUNT\_II\_SUPPORT}.
\hfill\textit{Sources:
\texttt{PatchedModuleAnnihilatorLe}\allowbreak\texttt{ComapIffIdealLe},
\texttt{AutomorphicDetection/LiftPt.core},
the original two \texttt{LinkTwo} inputs.}
% END PMS32

\subsection{Checked classical inputs and exact closure}

Only standard Noetherian commutative algebra is imported in
this chapter.  The primary-source certificate file gives the
precise tags and checked hypotheses for finite-module
completion, faithful flatness of local completion, Krull
intersection, finite-algebra decomposition over a complete
local ring, regular parameters and domains, height, associated
minimal primes, dimension of injective integral extensions,
and permanence of regularity and dimension under completion.
The determinant word identity is the already reconstructed
integral identity of JSP9.  The additional trace identity
needed for $\mathrm{ii}(W)\subset\ker\pi$ was proved in PMS30.

The new mathematical endpoints are the full original S5
(PMS15), the actual arithmetic PMod and all its prerequisites
on the supplied original data (PMS18--PMS29), the normalized
coefficient dimension (PMS16), and the precise kernel,
surjectivity and same-$J$ support comparison
(PMS30--PMS32).  The metadata distinguishes these completed
implications and producers from their unchanged original
TWAut, universal-ring and point inputs.
